\documentclass[11pt]{article}

\usepackage[T1]{fontenc}
\usepackage[margin=1in]{geometry}
\usepackage{amsmath,amssymb,mathtools,amsthm}
\usepackage{booktabs,tabularx,array}
\IfFileExists{lmodern.sty}
  {\usepackage{lmodern}\usepackage{microtype}}
  {\usepackage[expansion=false]{microtype}}
\usepackage{enumitem}
\usepackage[hidelinks]{hyperref}

\setlength{\parindent}{0pt}
\setlength{\parskip}{0.65em}
\setlist[itemize]{leftmargin=*,itemsep=0.2em,topsep=0.3em}
\setlist[enumerate]{leftmargin=*,itemsep=0.3em,topsep=0.3em}
\renewcommand{\arraystretch}{1.18}
\setcounter{tocdepth}{1}

\newcommand{\AP}{\operatorname{AP}}
\newcommand{\CNAP}{\operatorname{CNAP}}
\newcommand{\Regime}{\mathcal{R}}
\newcommand{\PSet}{\Pi_\star}
\newcommand{\E}{\mathbb{E}}
\newcommand{\Prb}{\mathbb{P}}
\newcommand{\iid}{\mathrel{\overset{\mathrm{iid}}{\sim}}}
\newcommand{\RobCNAP}{\underline{\CNAP}_{\PSet}}
\newcommand{\RobDelta}{\underline{\Delta}_{\PSet}}
\newcommand{\SRob}{S_{K,\PSet}^{\mathrm{AP}}(\Regime)}
\newcommand{\SRobHat}{\widehat S_{K,\PSet}^{\,\mathrm{AP}}}
\newcommand{\SRobNullHat}{\widehat S_{K,\PSet}^{\,\mathrm{AP,null}}}
\newcommand{\SRobCalHat}{\widehat S_{K,\PSet}^{\,\mathrm{AP,cal}}}

\newtheorem{definition}{Definition}
\newtheorem{proposition}{Proposition}
\theoremstyle{remark}
\newtheorem*{remark}{Remark}

\title{Supported Predictive Performance Across Target Prevalences\\
\large A Replication-Based Method for Average Precision}
\author{Marcus A. Thomas}
\date{Methodological proposal}

\begin{document}

\maketitle

\begin{abstract}
Predictive performance serves two tasks: interpreting the evidence for a model considered on its own and informing a decision between models evaluated together. Average precision makes both tasks sensitive to the prevalence of the population receiving the model. Prior standardization evaluates precision at each prevalence in a declared target set $\PSet$, and chance normalization creates a common CNAP scale with population random ranking at zero and perfect ranking at one. Each evaluation is summarized by its lowest CNAP over $\PSet$. For $K$ independent executions of an evaluation-generating regime $\Regime$, supported performance is
\[
S_{K,\PSet}^{\mathrm{AP}}(\Regime)
=
\E_{\Regime}
\left[
\min_{1\leq j\leq K}
\inf_{\pi\in\PSet}\CNAP_\pi(D_j)
\right].
\]
The reported level must survive every declared target prevalence and every replication. Widening $\PSet$ or increasing $K$ gives the claim additional opportunities for defeat. A stratified bootstrap estimates the quantity for fixed predictions on held-out outcomes. Conditional permutation locates one model relative to the chance level produced by the complete finite-design procedure, and the resulting calibrated value is the quantity reported for an isolated model. Compatible calibrated reports allow cautious benchmarking when joint predictions are unavailable. When two models are evaluated on the same units, supported superiority applies the operator to the lowest paired CNAP difference over $\PSet$, which stays on the raw CNAP scale because pairing already removes the shared design effect. A candidate meets a performance-based replacement criterion when its supported advantage exceeds a declared margin, with an outer lower confidence bound supplying population-level evidence. Pointwise prevalence profiles identify the populations that limit a claim and belong with the diagnostics. Two scalars carry the report: a calibrated robust value for the isolated model and a supported robust advantage for the choice between models.
\end{abstract}

\begingroup
\small
\setlength{\parskip}{0pt}
\tableofcontents
\endgroup

\begin{table}[ht]
\centering
\caption{Principal symbols.}
\label{tab:notation}
\begin{tabularx}{\textwidth}{@{}>{$}l<{$}X@{}}
\toprule
\text{Symbol} & Meaning \\
\midrule
\pi & A target-population prevalence \\
\PSet & The declared set of target prevalences that may challenge the claim \\
\Regime & The evaluation-generating regime, declaring what varies across replications \\
K & The number of independent replications that must retain the claim \\
\AP_\pi(D) & Prior-standardized average precision of evaluation $D$ at prevalence $\pi$ \\
\CNAP_\pi(D) & Chance-normalized $\AP_\pi$, zero at population random ranking and one at perfect ranking \\
\RobCNAP(D) & Lowest CNAP over $\PSet$, the prevalence-robust effect of one evaluation \\
\pi^\dagger(D) & A prevalence attaining that lowest value \\
Z_{\PSet} & The prevalence-robust effect of one replication drawn from $\Regime$ \\
S_K(T;\Regime) & Replication-survival operator, the mean of $\min_{j\leq K}T_j$ \\
\SRob & Prevalence-robust supported CNAP, the primary population quantity \\
\SRobHat & Its estimate from the bootstrap replication array \\
\SRobNullHat & The robust conditional chance level from permuted outcomes \\
\SRobCalHat & The calibrated value, reported for an isolated model \\
C_{\mathrm{prev}},\ C_{\mathrm{rep}} & The prevalence-challenge and replication-disagreement costs \\
\Delta_\pi(D) & Paired CNAP difference between models $A$ and $B$ at prevalence $\pi$ \\
\RobDelta(D) & Lowest paired difference over $\PSet$ \\
d & The smallest advantage that would justify replacing one model with another \\
\bottomrule
\end{tabularx}
\end{table}

\newpage

\section{Population conditions, corroboration, and model decisions}
\label{sec:motivation}

Many prediction systems rank cases so that a selected set can receive further attention. Candidate peptides may be selected for a vaccine, radiological images may be assigned for review, and customers may be prioritized for contact. Performance in these settings must describe the composition of the selected set as well as separation between outcome classes.

AUROC describes separation through the probability that a randomly drawn positive outranks a randomly drawn negative. It remains unchanged when prevalence changes and the class-conditional score distributions remain fixed. Selected-set purity behaves differently. Suppose a threshold selects $80\%$ of positive cases and $5\%$ of negative cases. In a balanced population of 10,000 cases, the selected set contains 4,000 positives and 250 negatives, giving precision $0.94$. In a population with $1\%$ positives, it contains 80 positives and 495 negatives, giving precision $0.14$. The threshold has the same true-positive and false-positive rates in both populations. The number of negative cases available for selection changes the meaning of its precision.

Average precision aggregates precision as the threshold is lowered. It therefore inherits the prevalence of the population being evaluated. Two studies can evaluate identical class-conditional ranking behavior and report different AP values because their populations contain different proportions of positive cases. Comparisons across studies require a common population interpretation.

Some applications have one well-defined target prevalence. Others serve several populations, settings, or planning scenarios whose prevalences differ. Selecting one convenient prevalence hides the populations in which a performance claim may weaken. Averaging over prevalences is the other obvious move, and it requires a distribution over target populations. Where such a distribution has a defensible scientific basis it should be used, and prevalence then belongs with the other sampled quantities discussed below. Where it does not, an assumed distribution lets strong performance in one population compensate for weak performance in another. A declared set of target prevalences makes the more demanding claim. The reported level must remain defensible throughout that set.

Replication supplies a second source of challenge. A fixed model can be tested on new evaluation units. A claim intended for new environments must also face environmental variation. A claim about a development procedure must survive new training, tuning, and model-selection runs. An evaluation-generating regime declares which of these mechanisms may vary. Independent executions of that regime give a performance claim repeated opportunities to fail.

These are two parts of one thing. A performance claim is stated against a declared \emph{challenge space}, the set of conditions under which it must survive. The target-prevalence set $\PSet$ declares the population conditions in that space. The regime $\Regime$ declares the data-generating and development conditions. The replication count $K$ declares how many independent executions must retain the claim. Together they give the resulting statement its scope, and a value quoted without them states nothing in particular.

The two parts are handled by different machinery, for a reason specific to average precision rather than a difference in kind. Section~\ref{sec:challenge-space} takes this up.

The available predictions determine how the statement enters a decision. One model evaluated on held-out outcomes can be located relative to the chance level of its finite evaluation design. Compatible reports from separate studies can then support cautious benchmarking. Competing models evaluated on the same units permit a paired analysis of the advantage retained throughout $\PSet$. The paired quantity supplies performance evidence for a choice or replacement decision.

\section{A common average-precision effect scale}
\label{sec:common-scale}

This section develops the first of the three declarations, the target-prevalence set $\PSet$. Sections~\ref{sec:regime} and \ref{sec:support} develop the regime $\Regime$ and the replication count $K$.

\subsection{Prior-standardized average precision}

At score threshold $c$, let
\[
r(c)=\Prb(S\geq c\mid Y=1),
\qquad
f(c)=\Prb(S\geq c\mid Y=0).
\]
These are the true-positive and false-positive rates. Precision in a population with prevalence $\pi\in(0,1)$ is
\begin{equation}
p_\pi(c)
=
\frac{\pi r(c)}
{\pi r(c)+(1-\pi)f(c)}.
\label{eq:reference-precision}
\end{equation}

Precision odds show how the population prior and ranking behavior enter:
\begin{equation}
\frac{p_\pi(c)}{1-p_\pi(c)}
=
\frac{\pi}{1-\pi}
\frac{r(c)}{f(c)}.
\label{eq:precision-odds}
\end{equation}
The first factor gives the target-population prior odds. The second is a threshold-specific likelihood ratio determined by the class-conditional score distributions. Prior standardization changes the first factor and preserves the second. The same factorization supports prior-probability and cost adjustment in classification \cite{elkan2001foundations}. It is also used to evaluate precision-based metrics at a fixed reference class ratio \cite{siblini2020calibration,degeest2016online}.

Applying Equation~\eqref{eq:reference-precision} across the distinct score thresholds gives
\[
\AP_\pi(D),
\]
the prior-standardized AP of evaluation $D$ at target prevalence $\pi$. The empirical definition and tie convention appear in Appendix~\ref{app:ap-definition}.

The transformation assumes that the class-conditional score distributions transport to the target populations represented by $\pi$. It describes prior-probability shift. Covariate shift and concept shift can change $r(c)$ and $f(c)$ and require a regime that represents those changes.

\subsection{Chance normalization}

Population random ranking selects the same fraction of each class, so $r(c)=f(c)$ and $p_\pi(c)=\pi$. Its prior-standardized AP is therefore $\pi$. Define
\begin{equation}
\CNAP_\pi(D)
=
\frac{\AP_\pi(D)-\pi}{1-\pi}.
\label{eq:cnap}
\end{equation}
CNAP is zero for population random ranking, one for perfect ranking, and negative below chance. Its range at prevalence $\pi$ is
\[
\left[-\frac{\pi}{1-\pi},1\right].
\]

The normalization has the general form
\[
\frac{\text{observed}-\text{chance}}
{\text{ceiling}-\text{chance}},
\]
which is familiar from chance-corrected agreement statistics \cite{cohen1960coefficient}. Earlier work has rescaled AP using either the observed prevalence or a fixed class ratio \cite{suyuan2015relationship,bestgen2021fisher}. Equation~\eqref{eq:cnap} first evaluates AP in the declared target population and then uses that population's random-ranking value as the lower anchor. Every target prevalence consequently shares the endpoints zero and one.

At empirical threshold $h$,
\begin{equation}
\frac{p_{\pi,h}-\pi}{1-\pi}
=
\frac{\pi(r_h-f_h)}
{\pi r_h+(1-\pi)f_h}.
\label{eq:threshold-cnap}
\end{equation}
Smaller prevalences give false positives greater influence because negative cases are more abundant. Larger prevalences credit class separation over a broader recall range. Model ordering can change with $\pi$ even after chance normalization.

\subsection{The target-prevalence challenge set}

Let
\[
\PSet\subset(0,1)
\]
be a nonempty compact set of target prevalences chosen from scientific knowledge, deployment plans, or a benchmark specification. A closed interval $[\pi_L,\pi_U]$ is a natural choice when every prevalence between two plausible limits is relevant. A finite set is appropriate when the application concerns named operating populations.

\begin{definition}[Prevalence-robust chance-normalized average precision]
\begin{equation}
\RobCNAP(D)
=
\inf_{\pi\in\PSet}\CNAP_\pi(D).
\label{eq:robust-cnap}
\end{equation}
\end{definition}

The quantity is the largest level retained at every prevalence in the declared set. The empirical CNAP profile is continuous on the compact set $\PSet$, so its infimum is attained. A minimizing prevalence may be recorded as
\[
\pi^\dagger(D)
\in
\arg\min_{\pi\in\PSet}\CNAP_\pi(D).
\]
Several minimizing values may exist.

When $\PSet=\{\pi_\star\}$, Equation~\eqref{eq:robust-cnap} reduces to the single-reference-prevalence effect. If $\PSet$ has maximum $\pi_U<1$, a common range is
\[
\RobCNAP(D)
\in
\left[-\frac{\pi_U}{1-\pi_U},1\right].
\]

The pointwise profile
\[
\pi\longmapsto\CNAP_\pi(D)
\]
shows which target populations limit the result and where model ordering changes. The profile serves as a diagnostic. Equation~\eqref{eq:robust-cnap} supplies the scalar claim.

The set must be declared independently of the observed performance profile. Its width controls the range of population conditions allowed to challenge the claim. If $\Pi_1\subseteq\Pi_2$, then
\[
\underline{\CNAP}_{\Pi_2}(D)
\leq
\underline{\CNAP}_{\Pi_1}(D).
\]
Widening the set cannot improve the retained level.

\subsection{Where the worst case occurs}
\label{sec:worst-case}

Appendix~\ref{app:ap-definition} shows that a threshold contribution is nondecreasing in $\pi$ whenever $r_h\geq f_h$. A model whose contributing thresholds all lie above the ROC diagonal therefore has a nondecreasing CNAP profile, and
\[
\inf_{\pi\in[\pi_L,\pi_U]}\CNAP_\pi(D)
=
\CNAP_{\pi_L}(D).
\]
Prevalence-robust performance then reduces to performance at the lower endpoint, and the declared set enters only through that endpoint.

Two consequences follow. For a single model, the robust scalar may carry little beyond a well-chosen point prevalence, and the choice of $\pi_L$ then requires the scientific justification that would otherwise attach to a single reference prevalence. For a paired comparison, the difference of two nondecreasing profiles need not be monotone. The limiting prevalence can fall in the interior and can differ across replications. The construction contributes most where model ordering changes across the target set.

How often interior minima occur is an empirical question. Recording the observed minimizer in each report allows it to be assessed across applications. Section~\ref{sec:paired} returns to this point, because the paired comparison is where a nonmonotone difference makes the declared set do work that no single prevalence could do.

\section{Performance across replications}
\label{sec:regime}

The second declaration is the regime. A performance claim acquires meaning from the conditions under which it could fail. These conditions are represented by an evaluation-generating regime.

\begin{definition}[Evaluation-generating regime]
An evaluation-generating regime $\Regime$ is a probability law for one complete replication of the claim under study. It declares which units, environments, measurement conditions, fitted models, and class counts are redrawn and which remain invariant.
\end{definition}

Three common regimes illustrate the range of claims.

\begin{table}[ht]
\centering
\caption{Illustrative evaluation-generating regimes.}
\label{tab:regimes}
\begin{tabularx}{\textwidth}{@{}>{\raggedright\arraybackslash}p{0.25\textwidth}XXX@{}}
\toprule
Component & Fixed-model & New-environment & Repeated-development \\
\midrule
Evaluation units & redrawn & redrawn & redrawn \\
Environment or period & invariant & redrawn & redrawn \\
Measurement process & invariant & declared & declared \\
Replication class counts & declared & declared & declared \\
Fitted model & invariant & invariant & redrawn \\
\bottomrule
\end{tabularx}
\end{table}

Every invariant component is a scientific conjecture. A fixed measurement process assumes that the relevant measurements behave comparably in the populations covered by the claim. A fixed fitted model confines the claim to that realized model. Repeated development expands the claim to the procedure that produced it. External replication can expose failures in these assumptions.

Write $n_+^{\mathrm{obs}}$ and $n_-^{\mathrm{obs}}$ for the observed evaluation counts. Let $n_+^{\mathrm{rep}}$ and $n_-^{\mathrm{rep}}$ denote the counts generated in one replication. A same-design analysis equates the replication counts with the observed counts. A prospective analysis declares the counts planned for future evaluations. The observed counts govern the information available for estimation. The replication counts help define the claim.

\subsection{Two parts of one challenge space}
\label{sec:challenge-space}

The set $\PSet$ is held fixed across replications. Every generated evaluation is examined at every prevalence in that set, prevalence being varied through Equation~\eqref{eq:reference-precision}. It is not drawn as an additional random variable.

This is not because prevalence is a different kind of challenge from the ones the regime carries. Both declare conditions under which the claim must hold. The difference is that prevalence is the only such condition that can be evaluated from a realized evaluation without generating a new one. The prior-odds factorization in Equation~\eqref{eq:precision-odds} means a single evaluation's threshold rates $(r_h,f_h)$ deliver $\CNAP_\pi$ at every $\pi$ at once. Nothing else in the challenge space behaves this way. What a model would have done in a different environment, under a different measurement process, or after a different development run cannot be computed from one environment's data. It has to be drawn.

The challenge space therefore has a computed part and a simulated part:
\[
\underbrace{\PSet}_{\text{evaluated by infimum}}
\qquad
\underbrace{\Regime,\ K}_{\text{explored by sampling}}.
\]
Worst case is taken over the whole space; expectation is taken over the sampled part alone. This is what Equation~\eqref{eq:robust-support} expresses, and it is why the infimum sits inside the replication rather than outside it.

The two remain separate in the notation, and $S_{K,\PSet}^{\mathrm{AP}}(\Regime)$ keeps them in distinct argument positions, because they are combined by different operators. Folding $\PSet$ into $\Regime$ would suggest averaging over prevalence, which is the construction Section~\ref{sec:motivation} declined.

\subsection{When separability fails}

Analytic separability is not free. It is the prior-probability-shift assumption of Section~\ref{sec:common-scale} under another description. Equation~\eqref{eq:reference-precision} reweights fixed class-conditional threshold rates, so it describes target populations that differ from the observed one in class proportions and in nothing else.

Where that assumption fails, prevalence does not leave the challenge space. It migrates from the computed part to the simulated part. If populations with different prevalences also differ in case mix, measurement, or score behavior, then $r(c)$ and $f(c)$ change with $\pi$, no reweighting recovers the target, and the environment component of the regime must generate the accompanying change directly. The declared claim is unaffected; only the machinery that interrogates it moves.

This gives a practical test for the boundary of the method. Ask whether the populations named by $\PSet$ plausibly differ in class proportions alone. Where they do, the prevalence challenge is computed and costs nothing. Where they do not, it must be simulated, and that requires data from the populations in question.

Let
\[
D^{\mathrm{rep}}\sim P_{\Regime}
\]
denote one execution of the regime and define
\begin{equation}
Z_{\PSet}
=
\RobCNAP(D^{\mathrm{rep}}).
\label{eq:robust-replication-effect}
\end{equation}
The law of $Z_{\PSet}$ is the prevalence-robust replication distribution. Its center describes expected worst-prevalence performance. Its shape describes how that retained level varies across replications.

Independence applies to complete replications. Observations within a replication may be dependent. Clustered, matched, longitudinal, and multi-environment designs require resampling at the exchangeable unit allowed by the design. A repeated-development regime reruns the complete fitting pipeline.

Replication-based assessment of classifier performance has precedent in work on generalization error and the replicability of learning experiments \cite{nadeau2003inference,bouckaert2004estimating}. The present replication distribution concerns a prevalence-robust CNAP effect generated under a declared regime.

\section{Replication-survival support}
\label{sec:support}

The third declaration is the replication count. This section builds the operator that consumes it, then combines all three declarations into the primary quantity.

Let $T$ be a scalar performance effect whose larger values indicate better performance. An evaluation yielding $t$ retains every claim of the form $T\geq s$ with $s\leq t$. Values above $t$ are defeated by that evaluation.

Suppose two independent replications yield $0.70$ and $0.05$. A claim based on both results cannot exceed $0.05$. The larger value supplies no compensation for a level defeated by the other replication. The greatest jointly retained level is their minimum.

For $K$ independent executions of $\Regime$, let
\[
T_1,\ldots,T_K
\iid
P_{\mathrm{rep}}^{\Regime,T}.
\]
Define
\begin{equation}
S_K(T;\Regime)
=
\E_{\Regime}
\left[
\min_{1\leq j\leq K}T_j
\right].
\label{eq:support-operator}
\end{equation}
The operator averages the greatest level retained by each set of $K$ replications. It remains on the scale of $T$.

This construction expresses corroboration through survival under independent opportunities for defeat \cite{popper1959logic,mayo1996error}. The regime determines the kinds of challenges presented to the claim. The value of $K$ determines how many independent challenges it faces. The default $K=2$ is the smallest value that requires joint survival.

If $T$ has lower bound $L$ and upper bound $U$, then
\begin{equation}
S_K(T;\Regime)
=
L+\int_L^U
\Prb_{\Regime}(T>s)^K\,ds.
\label{eq:support-survival}
\end{equation}
At each level $s$, the powered survival probability is the probability that all $K$ replications retain the claim. Appendix~\ref{app:proofs} gives the derivation.

At $K=2$, the minimum identity gives
\begin{equation}
S_2(T;\Regime)
=
\E_{\Regime}[T]
-
\frac12
\E_{\Regime}
\left[
|T_1-T_2|
\right].
\label{eq:k2-decomposition}
\end{equation}
Supported performance equals expected performance across replications minus half their expected absolute disagreement. The second term records the reduction in the retained level produced by variation across the regime.

\subsection{Prevalence-robust supported performance}

Apply Equation~\eqref{eq:support-operator} to the within-replication effect in Equation~\eqref{eq:robust-replication-effect}.

\begin{definition}[Prevalence-robust supported CNAP]
\begin{equation}
\SRob
=
\E_{\Regime}
\left[
\min_{1\leq j\leq K}
\inf_{\pi\in\PSet}
\CNAP_\pi(D_j)
\right].
\label{eq:robust-support}
\end{equation}
\end{definition}

The supported level survives every target prevalence in $\PSet$ and every one of the $K$ replications. The order of operations carries this interpretation. Each replication first receives the lowest value attained over the prevalence set. Replication support is then applied to those retained values.

Collapsing the two operations shows what the definition does to the challenge space of Section~\ref{sec:challenge-space}:
\[
\SRob
=
\E_{\Regime}
\left[
\inf_{\substack{\pi\in\PSet\\1\leq j\leq K}}
\CNAP_\pi(D_j)
\right].
\]
The infimum runs over the whole declared space, both the prevalences and the replications. The expectation runs over the sampled part alone. Read this way the definition is the only one available, and the ordering established in Section~\ref{sec:infimum-inside} is what happens when the two are taken in the wrong order.
If $L_{\PSet}$ is a common lower bound, its joint-survival form is
\begin{equation}
\SRob
=
L_{\PSet}
+
\int_{L_{\PSet}}^1
\Prb_{\Regime}
\left\{
\CNAP_\pi(D)>s
\text{ for every }\pi\in\PSet
\right\}^{K}
ds.
\label{eq:robust-joint-survival}
\end{equation}
The integrand gives the probability that every declared population condition retains level $s$ in all $K$ replications.

At $K=2$,
\begin{equation}
S_{2,\PSet}^{\mathrm{AP}}(\Regime)
=
\E_{\Regime}[Z_{\PSet}]
-
\frac12
\E_{\Regime}
\left[
|Z_{\PSet,1}-Z_{\PSet,2}|
\right].
\label{eq:robust-k2}
\end{equation}
The first term is expected worst-prevalence performance. The second records disagreement between the worst-prevalence levels attained in independent replications. Their minimizing prevalences may differ.

\subsection{Why the infimum belongs inside each replication}
\label{sec:infimum-inside}

Faced with a prevalence set and a replication criterion, the natural construction is to apply them in the order they were introduced. Calculate supported performance separately at each prevalence, giving the pointwise supported profile
\[
\pi
\longmapsto
S_K(\CNAP_\pi;\Regime),
\]
then take its lowest value,
\[
\inf_{\pi\in\PSet}S_K(\CNAP_\pi;\Regime).
\]
This is what a reader of a prevalence profile computes without thinking about it, and it is what a profile of separately supported values reports.

It is the weaker quantity. Every replication is challenged at the same prevalence before the prevalence minimum is taken, so a replication that fails badly at $\pi=0.01$ and a replication that fails badly at $\pi=0.40$ are never both charged for their own worst case. Only one prevalence is ever held against the whole array. Equation~\eqref{eq:robust-support} instead lets each realized replication reveal the target population that defeats it. The ordering is
\begin{equation}
S_K\left(
\inf_{\pi\in\PSet}\CNAP_\pi;
\Regime
\right)
\leq
\inf_{\pi\in\PSet}
S_K(\CNAP_\pi;\Regime).
\label{eq:strong-weak-order}
\end{equation}
The two sides agree when one prevalence minimizes the profile almost surely. They separate as the limiting prevalence moves between replications. Appendix~\ref{app:proofs} gives the short proof and shows that the entire gap comes from taking the expectation after the prevalence search rather than before.

The same reversal appears before any replication support is applied:
\begin{equation}
\E_{\Regime}
\left[
\inf_{\pi\in\PSet}\CNAP_\pi
\right]
\leq
\inf_{\pi\in\PSet}
\E_{\Regime}
\left[
\CNAP_\pi
\right].
\label{eq:prevalence-selection}
\end{equation}
Two named costs follow, and the paper uses these names throughout. The \emph{prevalence-challenge cost}
\[
C_{\mathrm{prev}}
=
\inf_{\pi\in\PSet}\E_{\Regime}[\CNAP_\pi]
-
\E_{\Regime}[Z_{\PSet}]
\]
is the difference between the two sides of Equation~\eqref{eq:prevalence-selection}. It is what the claim gives up by allowing the limiting prevalence to differ from one replication to the next, and it is zero when a single prevalence minimizes the realized profile almost surely. The \emph{replication-disagreement cost}
\[
C_{\mathrm{rep}}
=
\frac12
\E_{\Regime}
\left[
|Z_{\PSet,1}-Z_{\PSet,2}|
\right]
\]
is what the claim then gives up by requiring two replications to retain the same level, and it is the second term of Equation~\eqref{eq:robust-k2}. Together they give the reading used in Section~\ref{sec:reporting}:
\begin{equation}
S_{2,\PSet}^{\mathrm{AP}}(\Regime)
=
\underbrace{\inf_{\pi\in\PSet}\E_{\Regime}[\CNAP_\pi]}_{\text{anchor}}
-
C_{\mathrm{prev}}
-
C_{\mathrm{rep}}.
\label{eq:two-costs}
\end{equation}
The anchor is expected performance at the least favorable prevalence for the average profile. Both costs are nonnegative and both are available from the bootstrap array at no extra cost. Appendix~\ref{app:penalty} gives the derivation.

\subsection{Severity under prevalence and replication challenges}

The challenge space can be made more demanding along any of its three axes, and each move has a monotonicity result attached. If $\Pi_1\subseteq\Pi_2$, then
\begin{equation}
S_{K,\Pi_2}^{\mathrm{AP}}(\Regime)
\leq
S_{K,\Pi_1}^{\mathrm{AP}}(\Regime).
\label{eq:set-monotonicity}
\end{equation}
If $K_2\geq K_1$, then
\begin{equation}
S_{K_2,\PSet}^{\mathrm{AP}}(\Regime)
\leq
S_{K_1,\PSet}^{\mathrm{AP}}(\Regime).
\label{eq:k-monotonicity}
\end{equation}
Widening the prevalence set introduces more population conditions that can defeat the claim. Increasing $K$ introduces more independent replications that can defeat it. The third axis is the regime itself. A regime that spreads the replication distribution while holding its mean fixed also lowers the supported value, by the concave-order property in Appendix~\ref{app:families}. Declaring that measurement conditions or development runs may vary is therefore a severity choice of the same kind as widening $\PSet$ or raising $K$, though regimes are not ordered by inclusion and the comparison holds only between regimes ordered by dispersion.

Severity in each direction is a claim about the raw supported value. Section~\ref{sec:reporting} explains why the calibrated value cannot carry it.

The supported value does not give the probability that a claim is true. It is not a confidence bound and has no fixed recurrence frequency. It is an average retained performance level under declared challenges.

The replication-survival operator applies to any scalar effect oriented so that larger values indicate better performance. A loss $L(D)$ can be represented by $T(D)=-L(D)$. A paired loss comparison can use
\[
T(D)=L_B(D)-L_A(D),
\]
so positive values favor model $A$.

The prevalence challenge generalizes on the criterion of Section~\ref{sec:challenge-space} rather than on prevalence itself. Any dimension of a challenge space that can be evaluated from one realized evaluation, without generating another, can be handled by infimum at no sampling cost. Decision costs and case-mix weights are candidates wherever an analogue of Equation~\eqref{eq:precision-odds} isolates them in a single factor. Dimensions without that structure have to be simulated, and they belong in the regime.

\section{Estimation from one observed evaluation}
\label{sec:estimation}

The regime-indexed quantity in Equation~\eqref{eq:robust-support} is a population functional. One observed evaluation supplies a data-conditional representation of its replication distribution. The representation must reflect the mechanisms permitted to vary by the declared regime.

\subsection{Canonical fixed-model estimator}

Suppose $D_{\mathrm{obs}}$ is an independent or held-out evaluation of a fixed fitted model. Separate its positive and negative evaluation units into
\[
D_{\mathrm{obs},+},
\qquad
D_{\mathrm{obs},-},
\]
with counts $n_+^{\mathrm{obs}}>0$ and $n_-^{\mathrm{obs}}>0$. The canonical same-design estimator sets the replication counts equal to these observed counts.

For each bootstrap replication $b=1,\ldots,B$:
\begin{enumerate}
\item Sample $n_+^{\mathrm{rep}}$ units with replacement from $D_{\mathrm{obs},+}$.
\item Sample $n_-^{\mathrm{rep}}$ units with replacement from $D_{\mathrm{obs},-}$.
\item Combine the sampled units into $D^{(b)}$.
\item Evaluate the complete CNAP profile over $\PSet$.
\item Record
\begin{equation}
z_b
=
\inf_{\pi\in\PSet}
\CNAP_\pi(D^{(b)}).
\label{eq:bootstrap-robust-effect}
\end{equation}
\end{enumerate}

The array $z_1,\ldots,z_B$ is the empirical prevalence-robust replication distribution. Its construction holds the fitted model fixed and redraws evaluation units within outcome class. Hierarchical designs replace unit-level resampling with the exchangeable unit declared by the regime.

For $K=2$, the all-pairs estimator is
\begin{equation}
\widehat S_{2,\PSet,B}^{\,\mathrm{AP}}
=
\frac{2}{B(B-1)}
\sum_{1\leq i<j\leq B}
\min(z_i,z_j).
\label{eq:all-pairs-estimator}
\end{equation}
It is the order-two U-statistic with kernel $\min(z_i,z_j)$ \cite{hoeffding1948ustat}. It uses every unordered pair of independent computational replications. The equivalent sorted-array calculation appears in Appendix~\ref{app:computation}.

For general $K\leq B$, the corresponding U-statistic is
\begin{equation}
\widehat S_{K,\PSet,B}^{\,\mathrm{AP}}
=
\binom{B}{K}^{-1}
\sum_{\substack{I\subset\{1,\ldots,B\}\\|I|=K}}
\min_{i\in I}z_i.
\label{eq:general-u-statistic}
\end{equation}
An order-statistic formula avoids explicit enumeration.

The observed prevalence-robust effect is
\[
z_{\mathrm{obs}}
=
\inf_{\pi\in\PSet}
\CNAP_\pi(D_{\mathrm{obs}}).
\]
The observed minimizer $\pi^\dagger(D_{\mathrm{obs}})$ helps explain which target population limits the point estimate. The minimizing prevalence may vary across bootstrap replications, so it does not replace the full within-replication optimization.

\subsection{Continuous intervals and finite grids}

When $\PSet=[\pi_L,\pi_U]$, the empirical CNAP profile is continuous on a compact interval and attains its minimum. A one-dimensional optimizer can evaluate Equation~\eqref{eq:bootstrap-robust-effect}. The optimization tolerance should be small relative to the reported precision.

A prespecified finite grid
\[
\Pi_{\star,G}=\{\pi_1,\ldots,\pi_G\}
\]
defines a finite challenge set. Its minimum is exact for that set. If the grid approximates a continuous interval, report the grid and assess the approximation error. A coarse grid can miss an interior minimum and raise the reported robust value.

Every prevalence within a replication shares the same resampled units. This preserves the joint profile and allows the limiting prevalence to emerge from that replication. Independent resampling at separate prevalences would destroy this dependence.

\subsection{Prospective and repeated-development regimes}

A prospective fixed-model analysis replaces the observed bootstrap sizes with declared future evaluation counts and draws from the empirical class-conditional score pools. This changes the replication target and should be labeled as a prospective design.

Claims about new environments require data from multiple environments and a resampling hierarchy that redraws them. Claims about development procedures require complete reruns of fitting, tuning, and model selection. A bootstrap of one fixed score vector identifies neither target.

\section{Conditional calibration for one model}
\label{sec:calibration}

This section takes up the first of the two tasks, interpreting one model considered on its own. Section~\ref{sec:paired} takes up the second.

CNAP assigns zero to population random ranking. A finite evaluation design does not. Three mechanisms move its chance level away from zero, and they do not move it the same way. Stepwise AP has upward finite-sample bias under random ranking \cite{bestgen2015random}, which raises it. The replication-disagreement cost $C_{\mathrm{rep}}$ lowers it. The prevalence-challenge cost $C_{\mathrm{prev}}$ lowers it further. Only the complete pipeline determines where the three leave it, which is why the null is computed by running that pipeline rather than by correcting the value afterwards.

\subsection{A design with no signal at all}

The smallest case makes the point without any estimation. Take an evaluation with two observations, one positive and one negative, and let $\PSet=\{1/2\}$. There is no model worth speaking of and no association to find. Under random ranking the positive observation ranks first half the time and second half the time. Stepwise AP is then $1$ or $1/2$, so CNAP is $1$ or $0$.

The population random-ranking value is zero, but the replication mean is
\[
\E[Z_{\PSet}]
=
\frac12,
\]
and for two independent replications
\[
S_{2,\PSet}^{\mathrm{AP}}(\Regime_0)
=
\E[\min(Z_{\PSet,1},Z_{\PSet,2})]
=
\frac14.
\]
A procedure applied to pure noise returns a supported value of $0.25$ on a scale whose zero was supposed to mean chance. Nothing has gone wrong. The discreteness of stepwise AP puts half the mass at CNAP $=1$, and no amount of chance normalization can remove a distortion that lives in the finite design rather than in the population.

Two observations is an extreme case, and the effect shrinks as $n_+$ grows. It does not vanish, it interacts with ties and with $K$, and widening $\PSet$ adds the prevalence search to the same calculation. There is no general formula for where the three mechanisms leave the chance level of a particular design. The remedy is to run the whole procedure on outcomes that carry no signal and see what it returns.

\subsection{Conditional fixed-prediction null}

Suppose a developed model has been evaluated on outcomes that played no role in its development. Let $\mathbf{s}$ be its realized score vector and let $\mathbf{y}$ contain the observed positive and negative counts.

\begin{definition}[Conditional fixed-prediction null]
Condition on the score vector, its tie structure, the evaluation units, and the class counts. Assign outcomes to scores through the permutations allowed by the evaluation design.
\end{definition}

For independent evaluation units, every assignment preserving the class counts is allowed. Clustered and matched designs use restricted permutations at their exchangeable unit. The null asks how much prevalence-robust supported performance the procedure produces after the association between the realized scores and outcomes has been removed. Permutation tests for fixed classifier predictions provide the corresponding conditional no-association test \cite{ojala2010permutation}.

For each allowed or sampled permutation $m=1,\ldots,M$:
\begin{enumerate}
\item Attach the permuted outcomes to the fixed scores.
\item Construct $B$ bootstrap replications under the declared fixed-model regime.
\item Within every bootstrap replication, calculate
\[
z_b^{(m)}
=
\inf_{\pi\in\PSet}
\CNAP_\pi(D^{(m,b)}).
\]
\item Apply the support estimator to $z_1^{(m)},\ldots,z_B^{(m)}$ and record
\[
\widehat S_{K,\PSet,B}^{\,\mathrm{AP},(m)}.
\]
\end{enumerate}

The permutation replicates support two summaries, and the pair is easy to confuse.

The first is their mean, the estimated robust conditional chance level,
\begin{equation}
\SRobNullHat
=
\frac{1}{M}
\sum_{m=1}^{M}
\widehat S_{K,\PSet,B}^{\,\mathrm{AP},(m)}.
\label{eq:robust-null}
\end{equation}
This is the anchor that Section~\ref{sec:calibration-def} calibrates against.

The second is their upper tail,
\begin{equation}
p
=
\frac{
1+
\#\left\{
m:
\widehat S_{K,\PSet,B}^{\,\mathrm{AP},(m)}
\geq
\SRobHat
\right\}
}{M+1}.
\label{eq:permutation-value}
\end{equation}
This is a permutation $p$-value for the conditional no-association null. It is conditional on the realized score vector and its tie structure, so it refers to this evaluation rather than to any population, and it carries the usual cautions attaching to any $p$-value.

The two summaries use different features of the same $M$ replicates. Equation~\eqref{eq:robust-null} uses only their mean, so the calibrated value of Section~\ref{sec:calibration-def} discards the spread of the null distribution entirely. Equation~\eqref{eq:permutation-value} is the only quantity in the paper that uses that spread. Two designs with the same null mean and very different null dispersion produce identical calibrated values and quite different $p$, which is why both belong in a report.

The prevalence infimum, bootstrap support calculation, and permutation average form one procedure. Taking the infimum of separately calibrated pointwise values would use a different null construction and a different scale.

\subsection{Conditionally calibrated robust support}
\label{sec:calibration-def}

Perfect separation produces $\CNAP_\pi=1$ at every prevalence. Under the canonical stratified bootstrap, every bootstrap replication therefore has $z_b=1$ and the supported value is one.

\begin{definition}[Conditionally calibrated prevalence-robust supported CNAP]
\begin{equation}
\SRobCalHat
=
\frac{
\SRobHat-\SRobNullHat
}{
1-\SRobNullHat
}.
\label{eq:robust-calibration}
\end{equation}
\end{definition}

The calibrated score assigns the robust conditional chance level zero and perfect separation one. A value at or below the null anchor supplies no positive calibrated claim. Report $\SRobHat$ and $\SRobNullHat$ alongside the calibrated result, since a ratio cannot be checked without its numerator and denominator, and report the permutation $p$-value of Equation~\eqref{eq:permutation-value} for the dispersion the calibrated value omits.

Set inclusion has a direct severity interpretation for the raw quantity in Equation~\eqref{eq:robust-support}. Calibration changes both the observed supported value and its null anchor. The calibrated score can therefore move in either direction when $\PSet$ is widened. Statements about increased prevalence severity should refer to the raw supported value.

Two models can have different tie structures and different conditional null anchors. Their calibrated scales can consequently differ. Calibration locates one model relative to its finite-design chance level. A decision between models evaluated together uses their paired difference on the common uncalibrated CNAP scale.

The conditional null assumes that the outcomes used for calibration were isolated from model development. Reuse of evaluation outcomes for tuning, feature selection, stopping, or model choice invalidates the fixed-prediction permutation. A repeated-development claim requires a null procedure for the complete development process.

\section{Paired model comparison and replacement decisions}
\label{sec:paired}

The second task is a decision between two models evaluated together. It is not the first task performed twice. Calibration locates one model against the chance level of its own design, and two models evaluated on the same units have different tie structures and therefore different design chance levels. The comparison is made instead on the raw CNAP scale, where pairing removes the shared evaluation difficulty that calibration exists to handle for an isolated model.

Suppose models $A$ and $B$ are evaluated on the same units. At prevalence $\pi$, define the paired effect
\begin{equation}
\Delta_\pi(D)
=
\CNAP_{A,\pi}(D)
-
\CNAP_{B,\pi}(D).
\label{eq:paired-effect}
\end{equation}
Positive values favor model $A$. Shared changes in evaluation difficulty can cancel because both models are evaluated on the same realized data.

\begin{definition}[Prevalence-robust paired advantage]
\begin{equation}
\RobDelta(D)
=
\inf_{\pi\in\PSet}
\Delta_\pi(D).
\label{eq:robust-paired-effect}
\end{equation}
\end{definition}

The quantity is the largest advantage retained by model $A$ throughout the declared prevalence set. A minimizing value
\[
\pi_\Delta^\dagger(D)
\in
\arg\min_{\pi\in\PSet}\Delta_\pi(D)
\]
identifies a target population that limits the comparison.

Applying replication support gives
\begin{equation}
S_K(\RobDelta;\Regime)
=
\E_{\Regime}
\left[
\min_{1\leq j\leq K}
\inf_{\pi\in\PSet}
\Delta_\pi(D_j)
\right].
\label{eq:robust-paired-support}
\end{equation}
At $K=2$,
\begin{equation}
S_2(\RobDelta;\Regime)
=
\E_{\Regime}[\RobDelta]
-
\frac12
\E_{\Regime}
\left[
|\underline{\Delta}_{\PSet,1}-\underline{\Delta}_{\PSet,2}|
\right].
\label{eq:robust-paired-k2}
\end{equation}

For estimation, the same resampled units and outcomes are used for both models within bootstrap replication $b$. Calculate
\[
\delta_b
=
\inf_{\pi\in\PSet}
\left\{
\CNAP_{A,\pi}(D^{(b)})
-
\CNAP_{B,\pi}(D^{(b)})
\right\},
\]
then apply the all-pairs estimator to $\delta_1,\ldots,\delta_B$.

\subsection{Marginal supported values do not establish superiority}

Let
\[
A_{\PSet}(D)
=
\inf_{\pi\in\PSet}\CNAP_{A,\pi}(D),
\qquad
B_{\PSet}(D)
=
\inf_{\pi\in\PSet}\CNAP_{B,\pi}(D).
\]
For every evaluation,
\[
A_{\PSet}(D)-B_{\PSet}(D)
\geq
\RobDelta(D).
\]
At $K=2$, the ordering extends to
\begin{equation}
S_2(A_{\PSet};\Regime)
-
S_2(B_{\PSet};\Regime)
\geq
S_2(A_{\PSet}-B_{\PSet};\Regime)
\geq
S_2(\RobDelta;\Regime).
\label{eq:robust-marginal-order}
\end{equation}
Subtracting marginal supported values can therefore overstate the advantage retained under a paired comparison. Separate studies cannot recover the joint replication behavior needed for Equation~\eqref{eq:robust-paired-support}.

\subsection{Replacement criterion}

Let $d\geq0$ be the smallest retained CNAP advantage that would justify replacing model $B$ with model $A$. The performance criterion is
\begin{equation}
S_K(\RobDelta;\Regime)>d.
\label{eq:replacement}
\end{equation}
The choice $d=0$ requires a positive advantage throughout the prevalence set under the replication criterion. A positive $d$ states a practical threshold.

Some applications assign different practical meaning to the same CNAP difference at different prevalences. A prevalence-specific margin $d(\pi)$ can be incorporated through
\begin{equation}
S_K\left(
\inf_{\pi\in\PSet}
\{\Delta_\pi-d(\pi)\};
\Regime
\right)>0.
\label{eq:variable-margin}
\end{equation}
The margin function must be declared with the prevalence set.

The point estimate reports supported superiority under the empirical regime. Population-level evidence for replacement requires an outer lower confidence bound above $d$, or above zero for Equation~\eqref{eq:variable-margin}. Costs, feasibility, fairness, probability calibration, and operational consequences can also enter the final decision.

There is no general sharp permutation null for equality of two predictive models. Both models may carry signal and their score structures need not be exchangeable. Paired inference therefore targets Equation~\eqref{eq:robust-paired-support} directly.

\section{Interpretation, uncertainty, and reporting}
\label{sec:reporting}

\subsection{What the robust value carries}

The prevalence-robust supported value remains on the CNAP scale. It records a level retained across the population conditions in $\PSet$ and the replication challenges in $\Regime$. Its value depends on the replication counts, tie structure, class-conditional score distributions, and sources of variation represented by the regime.

Evaluation size enters through the replication distribution, and it enters through both costs at once. Section~\ref{sec:sparse} works through the case of few positive cases, where the effect is largest.

The robust value does not carry a fixed confidence level or recurrence probability. It does not protect against population changes outside $\PSet$. It also does not protect against changes in class-conditional score behavior unless the regime represents them.

\subsection{Sparse positive cases}
\label{sec:sparse}

Stepwise AP moves in increments of $1/n_+$. When positive cases are few, one rank change moves the profile at every prevalence, and three effects follow.

The replication distribution of $Z_{\PSet}$ becomes coarse with a heavy lower tail, so $C_{\mathrm{rep}}$ grows. The minimizing prevalence becomes unstable across replications, so $C_{\mathrm{prev}}$ grows as well. The conditional null rises, because the finite-sample upward bias of stepwise AP under random ranking is largest when $n_+$ is small.

Raw robust support therefore falls, while the calibrated value can move in either direction because its anchor also rises. Reporting $C_{\mathrm{prev}}$ and $C_{\mathrm{rep}}$ alongside the bootstrap distribution of $\pi^\dagger$ shows which mechanism is responsible. Neither a larger bootstrap count nor a narrower prevalence set recovers information the evaluation does not contain.

\subsection{Monte Carlo and estimation uncertainty}

The inner bootstrap approximates the data-conditional support functional. Its Monte Carlo error can be estimated from the U-statistic influence values in Appendix~\ref{app:computation}. The permutation average has a second Monte Carlo component governed by $M$.

An outer resampling layer addresses how well the observed evaluation identifies the supported target. Each outer sample must repeat the complete procedure:
\[
\text{outer resample}
\longrightarrow
\text{inner replication array}
\longrightarrow
\text{prevalence infima}
\longrightarrow
\text{support functional}.
\]
For calibrated performance, the robust conditional null must also be recomputed. For paired superiority, both models remain paired throughout the outer and inner layers.

The infimum can make the estimator nonsmooth when the minimizing prevalence changes. Sparse outcomes and flat minima can further affect interval behavior. Coverage should be studied by simulation under designs resembling the intended application.

\subsection{Reporting}

A report has three parts: what survived, what it was declared to survive, and what explains the number. The first is a sentence. The second is prespecified. The third belongs in the supplement. Producing every quantity in this paper for every analysis would obscure the claim it was built to make.

\subsubsection*{What survived}

For an isolated model the reported quantity is the calibrated value $\SRobCalHat$. The raw supported value and the conditional null travel with it and are never separated from it, because a ratio cannot be checked without its numerator and denominator. For a decision between two models the reported quantity is $S_K(\RobDelta;\Regime)$ on the raw CNAP scale. A single sentence carries either claim:

\begin{quote}
Across the prevalence range $[\pi_L,\pi_U]$ and $K$ independent replications of the declared regime, the model achieved a calibrated prevalence-robust CNAP of [value] (raw supported level [value] against a finite-design chance level of [value]). The worst case occurred at prevalence [value].
\end{quote}

The permutation $p$-value follows in the next sentence rather than inside this one. The quoted claim is a statement about magnitude, and a $p$-value placed within it invites the whole claim to be read as a significance result.

\subsubsection*{What it was declared to survive}

The prespecified declarations are $\PSet$ with the scientific basis for its boundaries, the AP convention and the prevalence-transport assumption, the regime with its class counts and exchangeable unit, $K$, the replacement margin $d$ where one applies, and the separation between model development and evaluation outcomes. These are design, not results. They are written once, before the analysis, and they belong in the methods rather than mixed among the estimates.

\subsubsection*{What explains the number}

Table~\ref{tab:reporting} places each remaining quantity against the question it answers.

\begin{table}[ht]
\centering
\caption{Which quantity answers which question.}
\label{tab:reporting}
\begin{tabularx}{\textwidth}{@{}XXl@{}}
\toprule
Question & Quantity & Where \\
\midrule
How much survived, on an interpretable scale? & $\SRobCalHat$ & headline \\
Could the null pipeline alone have produced this? & permutation $p$-value, Equation~\eqref{eq:permutation-value} & headline \\
Does $A$ beat $B$? & $S_K(\RobDelta;\Regime)$ & headline \\
Does the claim clear a threshold? & outer lower bound & only if a threshold exists \\
Why is the value low? & anchor, $C_{\mathrm{prev}}$, $C_{\mathrm{rep}}$, spread of $\pi^\dagger$ & supplement \\
Which populations limit it? & the CNAP profile & supplement \\
Did the computation converge? & Monte Carlo errors & methods \\
\bottomrule
\end{tabularx}
\end{table}

The outer layer is the item most often reported out of habit. It is required when a claim must clear a threshold, whether that threshold is a replacement margin $d$ over a competitor or an absolute minimum acceptable level for deployment, since $S_K(\RobDelta;\Regime)>d$ and $\SRob>c$ have the same shape and take the same treatment. Absent a threshold, a two-sided interval reports estimation uncertainty and a one-sided bound answers a question nobody asked. The distinction also has a practical edge. For paired superiority an outer resample costs one inner bootstrap, but for the calibrated value the conditional null must be recomputed inside every outer sample, so the cost carries a factor of $M$.

The profile should use the same prespecified set and the same shared replications as the scalar calculation.

\subsection{How to read the two headline values}

The raw supported value states how much performance survived. Use it to compare severity settings and to test a margin. The calibrated value states how far that level exceeds the mean chance level of this evaluation design. Use it to interpret one model. Both are magnitudes on a performance axis. Neither is a probability, and neither is a confidence bound.

The permutation $p$-value is a probability, and it answers the narrower question the calibrated value cannot. Calibration shifts and rescales by the null mean alone, so a supported value can sit above that mean and still fall well within the range the null pipeline produces. The $p$-value is what detects this. It is conditional on the realized score vector and its tie structure, so it speaks about this evaluation and not about the population. It is a subordinate diagnostic and not the claim.

Two cautions follow from Section~\ref{sec:calibration}. Widening $\PSet$ lowers the raw value monotonically, but the calibrated value can move either way because its anchor moves too, so a claim that a wider set was more severely tested must cite the raw value. And calibrated values are not differenced. Each is anchored to its own model's tie structure and design chance level, so two calibrated values compare only as an ordering. A difference between models requires the paired quantity of Section~\ref{sec:paired}.

The calibrated value is unavailable where the fixed-prediction permutation is invalid, which includes repeated-development regimes and any case where evaluation outcomes informed model development. The headline then reverts to the raw supported value, labelled as uncalibrated.

Separately reported calibrated values allow cautious benchmarking only when studies share $\PSet$, the AP convention, $K$, compatible regimes, comparable replication counts, and the same calibration procedure. Such reports remain marginal and data-conditional. They order models; they do not measure the gap between them.

\subsection{What an empirical assessment must establish}

The framework is developed here without an empirical illustration. Several questions must be settled before the robust scalar replaces a single-prevalence report in practice.

Matched evaluations with similar observed profiles should be compared to see whether their robust supported values separate. The conditional null must be located once the prevalence search is included, and its movement with the width of $\PSet$ recorded. Sparse-positive designs should be examined for their effect on $C_{\mathrm{prev}}$ and $C_{\mathrm{rep}}$. Paired superiority should be recomputed to see whether a conclusion drawn at one reference prevalence reverses. Outer intervals require a coverage study when the minimizer sits at an endpoint, in the interior, at several points, and when it moves under small perturbations. Monte Carlo stability requires the number of replications at which the supported value settles under prevalence optimization.

The last two items concern the estimator. The others determine whether the declared prevalence set changes any scientific conclusion. Section~\ref{sec:worst-case} identifies the case in which it cannot.

\section{Conclusion}

Average precision describes the purity of ranked selections in a target population. Its meaning changes with prevalence even when class-conditional ranking behavior remains fixed. Prior standardization evaluates that behavior throughout a declared set of target prevalences, and chance normalization gives every prevalence a common scale.

The prevalence-robust effect
\[
\inf_{\pi\in\PSet}\CNAP_\pi(D)
\]
records the level attained throughout the target set. Applying the replication-survival operator gives
\[
S_{K,\PSet}^{\mathrm{AP}}(\Regime)
=
\E_{\Regime}
\left[
\min_{1\leq j\leq K}
\inf_{\pi\in\PSet}
\CNAP_\pi(D_j)
\right].
\]
Each prevalence can challenge each replication. The supported value averages the maximal claim that survives all of them.

The three declarations of Section~\ref{sec:motivation} fix what that claim means. The prevalence set names the populations in which it must hold, the regime names the conditions under which it must hold, and the replication count says how many independent attempts at defeat it must survive. Widening the set or increasing the count makes the test more severe in a transparent direction. None of the three is estimated from the data, and a supported value quoted without them states nothing in particular.

The gain from this construction is uneven. When every contributing threshold of a single model lies above the ROC diagonal, its profile is nondecreasing and its worst case falls at the lower endpoint of an interval set. Paired differences can remain nonmonotone in the same evaluations. The robust comparison therefore retains content in cases where the robust single-model summary reduces to a point-prevalence report.

For one fixed model evaluated on held-out outcomes, stratified resampling estimates the prevalence-robust replication distribution. Conditional permutation measures the chance level produced by the complete finite-design procedure. The calibrated value locates the model relative to that reference. Compatible marginal reports can support cautious benchmarking when joint predictions are unavailable.

When both models are evaluated on the same units, the decision-relevant effect is
\[
\inf_{\pi\in\PSet}
\left\{
\CNAP_{A,\pi}(D)-\CNAP_{B,\pi}(D)
\right\}.
\]
Supported superiority applies the replication criterion to this paired effect. A replacement claim requires the supported advantage to exceed a declared margin, with an outer lower confidence bound supplying population-level evidence.

The two tasks therefore end on different scales, and for a reason. An isolated model is reported as a calibrated value because its raw level cannot be read without knowing where its own design puts chance. A comparison is reported on the raw paired scale because pairing has already removed the design effect that calibration was correcting. Prevalence profiles and the two costs explain those numbers when they are low, and belong with the diagnostics rather than in the claim.

What a study reports is one sentence:

\begin{quote}
This level survived every target prevalence declared relevant and every replication challenge represented by the regime.
\end{quote}

The value of the framework is that the sentence has a determinate meaning, and that a reader who disagrees with the declared populations or the declared regime can say so precisely.

\begin{thebibliography}{99}

\bibitem{bestgen2015random}
Bestgen, Y. (2015).
Exact expected average precision of the random baseline for system evaluation.
\emph{The Prague Bulletin of Mathematical Linguistics}, 103, 131--138.

\bibitem{bestgen2021fisher}
Bestgen, Y. (2021).
Improving measures of accuracy for detecting errors in texts.
\emph{Journal of Quantitative Linguistics}, 28, 293--310.

\bibitem{bouckaert2004estimating}
Bouckaert, R. R. (2004).
Estimating replicability of classifier learning experiments.
In \emph{Proceedings of the Twenty-First International Conference on Machine Learning}, 15.

\bibitem{cohen1960coefficient}
Cohen, J. (1960).
A coefficient of agreement for nominal scales.
\emph{Educational and Psychological Measurement}, 20, 37--46.

\bibitem{degeest2016online}
De Geest, R., Gavves, E., Ghodrati, A., Li, Z., Snoek, C., and Tuytelaars, T. (2016).
Online action detection.
In \emph{European Conference on Computer Vision}, 269--284.

\bibitem{elkan2001foundations}
Elkan, C. (2001).
The foundations of cost-sensitive learning.
In \emph{Proceedings of the Seventeenth International Joint Conference on Artificial Intelligence}, 973--978.

\bibitem{flach2015prg}
Flach, P. A. and Kull, M. (2015).
Precision-recall-gain curves: PR analysis done right.
In \emph{Advances in Neural Information Processing Systems}, 28.

\bibitem{hoeffding1948ustat}
Hoeffding, W. (1948).
A class of statistics with asymptotically normal distribution.
\emph{Annals of Mathematical Statistics}, 19, 293--325.

\bibitem{hosking1990lmoments}
Hosking, J. R. M. (1990).
L-moments: analysis and estimation of distributions using linear combinations of order statistics.
\emph{Journal of the Royal Statistical Society, Series B}, 52, 105--124.

\bibitem{mayo1996error}
Mayo, D. G. (1996).
\emph{Error and the Growth of Experimental Knowledge}.
University of Chicago Press.

\bibitem{nadeau2003inference}
Nadeau, C. and Bengio, Y. (2003).
Inference for the generalization error.
\emph{Machine Learning}, 52, 239--281.

\bibitem{ojala2010permutation}
Ojala, M. and Garriga, G. C. (2010).
Permutation tests for studying classifier performance.
\emph{Journal of Machine Learning Research}, 11, 1833--1863.

\bibitem{popper1959logic}
Popper, K. R. (1959).
\emph{The Logic of Scientific Discovery}.
Hutchinson.

\bibitem{siblini2020calibration}
Siblini, W., Fr\'ery, J., He-Guelton, L., Obl\'e, F., and Wang, Y. (2020).
Master your metrics with calibration.
In \emph{Advances in Intelligent Data Analysis XVIII}, 457--469.

\bibitem{suyuan2015relationship}
Su, W., Yuan, Y., and Zhu, M. (2015).
A relationship between the average precision and the area under the ROC curve.
In \emph{Proceedings of the 2015 International Conference on The Theory of Information Retrieval}, 349--352.

\bibitem{wang1996premium}
Wang, S. (1996).
Premium calculation by transforming the layer premium density.
\emph{ASTIN Bulletin}, 26, 71--92.

\end{thebibliography}

\appendix

\section{Prior-standardized average precision}
\label{app:ap-definition}

\subsection{Empirical definition}

Let
\[
D=\{(s_i,y_i)\}_{i=1}^{n},
\qquad
y_i\in\{0,1\},
\]
with $n_+=\sum_i y_i>0$ and $n_-=n-n_+>0$. Let
\[
c_1>c_2>\cdots>c_H
\]
be the distinct observed score values. Threshold $h$ selects every observation with $s_i\geq c_h$. Define
\[
\operatorname{TP}_h
=
\sum_{i=1}^{n}
\mathbf{1}\{y_i=1,\ s_i\geq c_h\},
\qquad
\operatorname{FP}_h
=
\sum_{i=1}^{n}
\mathbf{1}\{y_i=0,\ s_i\geq c_h\},
\]
and
\[
r_h=\frac{\operatorname{TP}_h}{n_+},
\qquad
f_h=\frac{\operatorname{FP}_h}{n_-}.
\]
Prior-standardized precision at threshold $h$ is
\[
p_{\pi,h}
=
\frac{\pi r_h}
{\pi r_h+(1-\pi)f_h}.
\]
Set $r_0=0$. The non-interpolated prior-standardized AP is
\begin{equation}
\AP_\pi(D)
=
\sum_{h=1}^{H}
(r_h-r_{h-1})p_{\pi,h}.
\label{eq:empirical-ap}
\end{equation}
Equal scores enter as one threshold block. A threshold that adds only negative cases has zero recall increment. Its false positives affect later precision values.

Using $H$ for the threshold count avoids conflict with the replication severity parameter $K$.

\subsection{Weighted implementation}

The same quantity can be calculated by weighting the observed classes to represent prevalence $\pi$. Assign each positive observation weight
\[
w_+=\frac{\pi}{n_+}
\]
and each negative observation weight
\[
w_-=\frac{1-\pi}{n_-}.
\]
At threshold $h$, weighted precision is
\[
\frac{w_+\operatorname{TP}_h}
{w_+\operatorname{TP}_h+w_-\operatorname{FP}_h}
=
\frac{\pi r_h}
{\pi r_h+(1-\pi)f_h}.
\]
Weighted recall remains $r_h$, so stepwise weighted AP equals Equation~\eqref{eq:empirical-ap}.

\subsection{Behavior across target prevalences}

Combining Equations~\eqref{eq:cnap} and \eqref{eq:empirical-ap} gives
\begin{equation}
\CNAP_\pi(D)
=
\sum_{h=1}^{H}
(r_h-r_{h-1})
\frac{\pi(r_h-f_h)}
{\pi r_h+(1-\pi)f_h}.
\label{eq:empirical-cnap-profile}
\end{equation}
The derivative of one threshold contribution is
\begin{equation}
\frac{\partial}{\partial\pi}
\left\{
\frac{\pi(r_h-f_h)}
{\pi r_h+(1-\pi)f_h}
\right\}
=
\frac{f_h(r_h-f_h)}
{\{\pi r_h+(1-\pi)f_h\}^2}.
\label{eq:profile-derivative}
\end{equation}
A threshold with $r_h\geq f_h$ has a nondecreasing contribution. If this holds at every threshold with a positive recall increment, the complete CNAP profile is nondecreasing and its minimum over $[\pi_L,\pi_U]$ occurs at $\pi_L$. Profiles can have interior extrema when contributing thresholds lie on different sides of the ROC diagonal.

Two nondecreasing marginal profiles can have a nonmonotone difference. Prevalence-robust paired comparison can therefore remain informative when each single-model infimum occurs at the lower endpoint.

\section{Related precision standardizations}
\label{app:related}

Equation~\eqref{eq:precision-odds} isolates prevalence in the prior-odds factor. Prior standardization assigns that factor a declared value and retains the selected-set interpretation of precision. Precision-recall-gain analysis removes the factor from its threshold-level precision coordinate \cite{flach2015prg}.

Let
\[
p
=
\frac{\pi r}{\pi r+(1-\pi)f}.
\]
Precision gain is
\[
\operatorname{precG}
=
\frac{p-\pi}{(1-\pi)p}.
\]
Writing $D=\pi r+(1-\pi)f$ gives
\[
p-\pi
=
\frac{\pi(1-\pi)(r-f)}{D},
\qquad
(1-\pi)p
=
\frac{(1-\pi)\pi r}{D},
\]
and hence
\begin{equation}
\operatorname{precG}
=
\frac{r-f}{r}
=
1-\frac{f}{r}.
\label{eq:precision-gain}
\end{equation}
The prevalence terms cancel at the threshold level.

For $r>0$, the CNAP threshold contribution in Equation~\eqref{eq:threshold-cnap} satisfies
\begin{equation}
\lim_{\pi\to1}
\frac{\pi(r-f)}
{\pi r+(1-\pi)f}
=
1-\frac{f}{r}.
\label{eq:precision-gain-limit}
\end{equation}
Precision gain is therefore the high-prevalence limit of the threshold-level normalized precision used by CNAP.

The integral summaries remain different. CNAP averages its threshold contribution against recall increments. A precision-recall-gain area uses recall gain and its associated integration measure. The two constructions can consequently order models differently. The present framework retains target prevalence because selected-set purity in the populations represented by $\PSet$ carries the scientific interpretation of interest.

\section{Properties of prevalence-robust support}
\label{app:proofs}

\subsection{Joint-survival representation}

For a random variable $T\in[L,U]$,
\[
T
=
L+\int_L^U\mathbf{1}\{T>s\}\,ds.
\]
For independent copies $T_1,\ldots,T_K$,
\[
\min_jT_j
=
L+
\int_L^U
\prod_{j=1}^{K}
\mathbf{1}\{T_j>s\}\,ds.
\]
Taking expectations and using independence gives Equation~\eqref{eq:support-survival}.

Set $T=\inf_{\pi\in\PSet}\CNAP_\pi(D)$. The event $T>s$ is equivalent to
\[
\CNAP_\pi(D)>s
\quad
\text{for every }\pi\in\PSet.
\]
Substitution gives Equation~\eqref{eq:robust-joint-survival}.

\subsection{Set monotonicity}

If $\Pi_1\subseteq\Pi_2$, then for every evaluation
\[
\inf_{\pi\in\Pi_2}\CNAP_\pi(D)
\leq
\inf_{\pi\in\Pi_1}\CNAP_\pi(D).
\]
The minimum over $K$ replications preserves the inequality, as does expectation. This proves Equation~\eqref{eq:set-monotonicity}.

\subsection{Ordering relative to the supported profile}

Let $X_\pi^{(j)}=\CNAP_\pi(D_j)$. For every realized set of replications,
\[
\min_j\inf_{\pi\in\PSet}X_\pi^{(j)}
=
\inf_{\substack{\pi\in\PSet\\1\leq j\leq K}}
X_\pi^{(j)}
=
\inf_{\pi\in\PSet}\min_jX_\pi^{(j)}.
\]
The two orderings of the minimum agree pathwise, so no loss occurs at this step. Write
\[
Y_\pi=\min_{1\leq j\leq K}X_\pi^{(j)}
\]
and take expectations. Since $\inf_\pi Y_\pi\leq Y_{\pi'}$ for every $\pi'\in\PSet$,
\[
\E\left[\inf_{\pi\in\PSet}Y_\pi\right]
\leq
\inf_{\pi\in\PSet}\E[Y_\pi],
\]
which proves Equation~\eqref{eq:strong-weak-order}. The gap arises entirely at this step. It is the cost of taking the expectation after the prevalence search rather than before.

\subsection{The $K=2$ decomposition}

For real values $a$ and $b$,
\[
\min(a,b)
=
\frac{a+b-|a-b|}{2}.
\]
Apply this identity to two independent copies of $Z_{\PSet}$ and take expectations to obtain Equation~\eqref{eq:robust-k2}.

\subsection{Two costs of the prevalence challenge}
\label{app:penalty}

Equation~\eqref{eq:two-costs} follows by adding and subtracting $\E_{\Regime}[Z_{\PSet}]$ in Equation~\eqref{eq:robust-k2}. The prevalence-challenge cost $C_{\mathrm{prev}}$ is nonnegative by Equation~\eqref{eq:prevalence-selection}. The replication-disagreement cost $C_{\mathrm{rep}}$ is nonnegative because it is half a mean absolute difference.

Estimation uses the same bootstrap array that produces $\SRobHat$. The anchor requires the mean CNAP profile across replications, evaluated on the shared prevalence grid, together with its minimizer. Note that this minimizer is a property of the averaged profile and generally differs from any individual $\pi^\dagger(D^{(b)})$. The two costs then follow by subtraction, since $\E[Z_{\PSet}]$ is the mean of the array and $S_{2,\PSet}$ is its all-pairs statistic.

Separating them answers a question the supported value alone cannot. A result limited by $C_{\mathrm{prev}}$ is limited by an unstable worst-case population, which more evaluation units may not repair if the profile is genuinely flat near its minimum. A result limited by $C_{\mathrm{rep}}$ is limited by replication variability, which a larger evaluation usually does repair under a fixed-model regime. The two call for different responses.

\subsection{Affine equivariance}

For $a>0$ and constant $c$,
\[
\inf_{\pi\in\PSet}\{aX_\pi+c\}
=
a\inf_{\pi\in\PSet}X_\pi+c.
\]
The support operator also satisfies
\[
S_K(aT+c;\Regime)
=
aS_K(T;\Regime)+c.
\]
These identities explain why the common pointwise chance normalization is compatible with the prevalence infimum and replication support. A prevalence-dependent affine transformation would not commute with the infimum.

\subsection{Two established families}
\label{app:families}

Let $Q$ be the quantile function of $Z_{\PSet}$. The minimum of $K$ independent copies equals $Q(U_{(1)})$, where $U_{(1)}$ is the smallest of $K$ independent uniform variables and has density $K(1-v)^{K-1}$. Hence
\begin{equation}
S_{K,\PSet}^{\mathrm{AP}}(\Regime)
=
\int_0^1
Q(v)\,K(1-v)^{K-1}\,dv.
\label{eq:quantile-weight}
\end{equation}
The weight decreases in $v$, so the operator places its mass on the lower tail of the replication distribution.

Two identifications follow and supply properties without separate proof.

\paragraph{L-moments.}
The first two L-moments of a distribution \cite{hosking1990lmoments} are $\lambda_1=\E[Z_{\PSet}]$ and $\lambda_2=\tfrac12\E|Z_{\PSet,1}-Z_{\PSet,2}|$, the latter being half the Gini mean difference. Equation~\eqref{eq:robust-k2} is therefore the exact identity
\[
S_{2,\PSet}^{\mathrm{AP}}(\Regime)
=
\lambda_1-\lambda_2 .
\]
L-moment estimators are linear in order statistics, which is the origin of the sorted-sum form in Appendix~\ref{app:computation}, and they exist whenever the first moment does.

\paragraph{Distortion measures.}
Equation~\eqref{eq:support-survival} applies the distortion $u\mapsto u^{K}$ to the survival function of $Z_{\PSet}$, placing $S_K$ in the dual-power family of Wang \cite{wang1996premium}. The following properties hold.

\begin{itemize}
\item \emph{Monotonicity under first-order stochastic dominance.} If $Q_A(v)\geq Q_B(v)$ for all $v\in(0,1)$, then Equation~\eqref{eq:quantile-weight} gives $S_K(A;\Regime)\geq S_K(B;\Regime)$. Section~\ref{sec:paired} uses this property.
\item \emph{Monotonicity under the concave order.} The quantile weight is nonincreasing, so a mean-preserving spread of the replication distribution cannot raise the supported value.
\item \emph{Translation equivariance and positive homogeneity.} These restate the affine identities above.
\item \emph{Comonotone additivity.} If $X$ and $Y$ are comonotone, then $Q_{X+Y}=Q_X+Q_Y$ and $S_K(X+Y)=S_K(X)+S_K(Y)$.
\item \emph{Superadditivity.} For any $X$ and $Y$ evaluated on the same replications, $\min_j(X_j+Y_j)\geq\min_jX_j+\min_jY_j$, so
\[
S_K(X+Y;\Regime)
\geq
S_K(X;\Regime)+S_K(Y;\Regime).
\]
Section~\ref{sec:paired} uses this property with $X=A_{\PSet}-B_{\PSet}$ and $Y=B_{\PSet}$.
\end{itemize}

The replication-survival criterion determines the operator. These representations place the resulting functional in familiar families and supply the properties used elsewhere in the paper.

\subsection{Marginal and paired robust support}

Let $A_{\PSet}$, $B_{\PSet}$, and $\RobDelta$ be defined as in Section~\ref{sec:paired}. Choose any prevalence minimizing $A_\pi$. Then
\[
\RobDelta
\leq
A_{\PSet}-B_\pi
\leq
A_{\PSet}-B_{\PSet}.
\]
Thus $A_{\PSet}-B_{\PSet}\geq\RobDelta$ pointwise, and monotonicity of $S_2$ gives
\[
S_2(A_{\PSet}-B_{\PSet};\Regime)
\geq
S_2(\RobDelta;\Regime).
\]

For any integrable random variables $A$ and $B$,
\[
S_2(A)-S_2(B)
=
\E[A-B]
-
\frac12
\left\{
\E|A_1-A_2|
-
\E|B_1-B_2|
\right\}.
\]
The triangle inequality gives
\[
|A_1-A_2|
\leq
|(A_1-B_1)-(A_2-B_2)|
+
|B_1-B_2|.
\]
Rearrangement yields
\[
S_2(A)-S_2(B)
\geq
S_2(A-B).
\]
Taking $A=A_{\PSet}$ and $B=B_{\PSet}$ proves Equation~\eqref{eq:robust-marginal-order}.

\subsection{Relation to a standard-error adjustment}

Suppose $Z_{\PSet}$ is normally distributed with mean $\mu$ and standard deviation $\sigma$. Two independent copies have difference
\[
Z_{\PSet,1}-Z_{\PSet,2}
\sim
N(0,2\sigma^2),
\]
so
\[
\frac12
\E
\left[
|Z_{\PSet,1}-Z_{\PSet,2}|
\right]
=
\frac{\sigma}{\sqrt{\pi}}.
\]
Equation~\eqref{eq:robust-k2} becomes
\[
S_{2,\PSet}^{\mathrm{AP}}(\Regime)
=
\mu-\frac{\sigma}{\sqrt{\pi}}.
\]
This numerical relation belongs to the normal family. The general replication-disagreement cost depends on the complete replication distribution and is not a universal multiple of its standard deviation.

\subsection{The support distribution and frequency-qualified levels}

Let
\[
M_K
=
\min\{Z_{\PSet,1},\ldots,Z_{\PSet,K}\}.
\]
Its distribution records the level jointly retained by $K$ replications. Prevalence-robust supported performance is $\E[M_K]$.

A frequency-qualified retained level can be defined by
\[
m_\gamma^{(K)}
=
\sup
\left\{
m:
\Prb(M_K\geq m)\geq\gamma
\right\}.
\]
If $Q_{\PSet}$ is the quantile function of $Z_{\PSet}$ and the distribution is continuous, then
\begin{equation}
m_\gamma^{(K)}
=
Q_{\PSet}\left(1-\gamma^{1/K}\right).
\label{eq:frequency-level}
\end{equation}
This quantity answers a different question from the mean retained level. It reports a level attained jointly with frequency $\gamma$ under the regime. It remains data-conditional when estimated from a bootstrap replication array and is not a confidence bound for a population parameter.

\subsection{Contraction of the replication distribution}

Suppose a sequence of evaluation designs makes $Z_{\PSet}$ converge in probability to a stable value $z_0$ and the effects remain uniformly bounded. Then
\[
S_{K,\PSet}^{\mathrm{AP}}(\Regime)
\longrightarrow
z_0.
\]
Larger evaluations can produce this contraction when the regime redraws independent units from stable class-conditional score distributions. Environmental variation, clustered dependence, and repeated model development can prevent it. Replication counts therefore belong in the regime description even though the supported effect contains no explicit sample-size multiplier.

\section{Regime construction}
\label{app:regime}

Let $E$ index environments and $U$ index evaluation units within an environment. A regime can be written schematically as
\[
E_j\sim P_E,
\qquad
U_{ij}\mid E_j\sim P_{U\mid E_j},
\qquad
D^{\mathrm{rep}}=\{(S_{ij},Y_{ij})\}.
\]
A fixed-environment analysis conditions on the observed $E$. A new-environment analysis draws $E_j$ from a population represented by the available environments. Cluster resampling preserves the dependence of units within each environment.

The target-prevalence set remains in the computed part of the challenge space throughout this construction. Each generated dataset produces class-conditional threshold rates, and Equation~\eqref{eq:reference-precision} evaluates those rates throughout $\PSet$ without further sampling.

Section~\ref{sec:challenge-space} gives the condition under which this is legitimate and the remedy when it is not. In the notation above, migration means that $\pi$ ceases to be a free argument of the evaluation and becomes a property of $E_j$. The environment law $P_E$ then has to place mass on environments whose prevalences span $\PSet$, and the analysis requires data from those environments. A regime that redraws units within a single environment cannot recover this, whatever range of $\pi$ is declared.

Repeated-development regimes add training data, fitting randomness, tuning, and selection:
\[
D_{\mathrm{train}}^{(j)}
\longrightarrow
\widehat f^{(j)}
\longrightarrow
D_{\mathrm{eval}}^{(j)}
\longrightarrow
\inf_{\pi\in\PSet}
\CNAP_\pi(\widehat f^{(j)},D_{\mathrm{eval}}^{(j)}).
\]
Each arrow belongs to the replication. Holding any component fixed narrows the resulting claim.

\section{Efficient computation}
\label{app:computation}

\subsection{All-pairs support for $K=2$}

Sort the prevalence-robust replication effects:
\[
z_{(1)}\leq\cdots\leq z_{(B)}.
\]
The value $z_{(i)}$ is the minimum in each pair with a larger-indexed value, so
\begin{equation}
\widehat S_{2,\PSet,B}^{\,\mathrm{AP}}
=
\frac{2}{B(B-1)}
\sum_{i=1}^{B-1}
(B-i)z_{(i)}.
\label{eq:sorted-pairs}
\end{equation}
Sorting costs $O(B\log B)$ and the weighted sum costs $O(B)$.

\subsection{General $K$}

For the order-$K$ U-statistic, $z_{(i)}$ is the minimum for every subset containing it and choosing the remaining $K-1$ values from the $B-i$ larger values. Hence
\begin{equation}
\widehat S_{K,\PSet,B}^{\,\mathrm{AP}}
=
\binom{B}{K}^{-1}
\sum_{i=1}^{B-K+1}
\binom{B-i}{K-1}
z_{(i)}.
\label{eq:sorted-general-k}
\end{equation}

\subsection{Monte Carlo standard error}

For $K=2$, define the leave-one replication projection
\[
\widehat h_{1,-i}(z_i)
=
\frac{1}{B-1}
\sum_{j\ne i}
\min(z_i,z_j).
\]
A first-order Monte Carlo standard error is
\begin{equation}
\widehat{\operatorname{se}}_{\mathrm{MC}}
=
\frac{2}{\sqrt B}
\operatorname{sd}_{1\leq i\leq B}
\left\{
\widehat h_{1,-i}(z_i)
\right\}.
\label{eq:mcse}
\end{equation}
Prefix sums over the sorted array evaluate all projections in linear time after sorting.

\subsection{Prevalence evaluation}

For a finite set with $G$ values, evaluating every prevalence within every bootstrap replication costs $O(BGH)$ after threshold counts are available, where $H$ is the number of distinct score thresholds. The values share the same $r_h$ and $f_h$, so vectorization over $\pi$ avoids rebuilding the ranking.

For a continuous interval, Equation~\eqref{eq:empirical-cnap-profile} supplies a smooth one-dimensional objective away from irrelevant zero-recall thresholds. Endpoint values and stationary points can be checked directly. A numerical optimizer should be verified against a dense grid during implementation.

Conditional calibration repeats the robust support procedure for each permutation. The permutations are independent computational jobs. Monte Carlo stability of the null mean usually requires fewer permutations than accurate estimation of a small tail probability.

\section{Calibration and interval considerations}
\label{app:calibration}

\subsection{Perfect-separation fixed point}

If every positive score exceeds every negative score, every bootstrap sample preserves that ordering. At each target prevalence,
\[
\AP_\pi=1,
\qquad
\CNAP_\pi=1.
\]
The prevalence infimum and every replication minimum equal one. The canonical raw and calibrated supported values therefore attain their upper anchor.

\subsection{Conditional validity}

The fixed-prediction permutation distribution is valid when the evaluation outcomes are exchangeable with respect to the fixed scores under the null and the permutation respects the design. The procedure conditions on the realized score vector, class counts, and tie structure. It does not account for outcome information used during model development.

The robust null includes the search over $\PSet$. Every permutation must use the same prevalence set, numerical optimization, resampling counts, and support severity as the observed calculation. Altering the set or optimization after inspecting the observed profile breaks the correspondence.

\subsection{Outer intervals}

An outer bootstrap sample yields a new empirical distribution of class-conditional scores. The complete inner estimator produces one value of the target statistic for that outer sample. Percentile intervals are simple. Basic, studentized, and bias-corrected intervals may respond differently to skewness and changes in the minimizing prevalence.

For paired superiority, the outer resampling preserves the pairing between model scores. A lower interval endpoint above the margin supports the replacement criterion at the interval's stated confidence level. Simulation should examine coverage when the prevalence minimum occurs at an endpoint, in the interior, at several points, or changes under small perturbations.

\section{Reporting templates}
\label{app:report-template}

\subsection{Single-model performance}

A concise report can follow this sequence:
\begin{quote}
The model was evaluated over the prespecified target-prevalence set $\PSet$. Its observed prevalence-robust CNAP was [value], with a minimizing prevalence of [value or set]. Under the declared [same-design or prospective] regime with replication counts [counts] and $K=[K]$, prevalence-robust supported CNAP was [value]. The conditional permutation null was [value], giving calibrated robust support [value]. The permutation $p$-value was [value] and the outer [level]\% interval was [limits].
\end{quote}
The accompanying methods identify the AP convention, resampling unit, optimization method, bootstrap and permutation counts, Monte Carlo errors, and separation of model development from evaluation.

\subsection{Paired replacement decision}

A paired report can follow this sequence:
\begin{quote}
Model $A$ was compared with model $B$ on the same evaluation units over $\PSet$. The observed worst-prevalence CNAP advantage was [value], attained at [prevalence]. Under the declared regime and $K=[K]$, supported superiority was [value]. The prespecified replacement margin was [value], and the outer [level]\% interval was [limits].
\end{quote}
The report states whether the lower interval endpoint exceeds the margin and describes the additional considerations entering the decision.

\end{document}
